\section{Conclusion}

This paper reframes operational demand planning as feasibility-constrained model selection through a forecast-to-plan evaluation gate. The central finding is that in 91.1\% of evaluation groups, the accuracy-best deployable strategy diverges from the operational-loss-optimal deployable strategy. The accuracy-first execution violation rate (15.4\% of periods) is 5.7 times higher than the best deployable operational-loss strategy (2.7\% of periods).

Section~\ref{sec:gap} opened with four operational symptoms, and each now has a direct counterpart above rather than a generic claim of improvement. \emph{Execution violations} are what the PEG-rate metric catches (15.4\% versus 2.7\%). \emph{Governance burden} is addressed by the protocol's governance-maturity lane, which recommends DP or budgeted DP specifically where switching requires approval, on the empirical basis of the greedy-to-DP gap. \emph{Supplier disruption} falls where planning volatility falls---Simple Ensemble cuts execution violations by roughly 37 percentage points at no WAPE cost, though other cases require a small WAPE cost for stability, which the execution-flexibility lane makes explicit. \emph{Operational gridlock} is addressed structurally: the five-step protocol replaces ad hoc negotiation with a repeatable, auditable procedure. Each symptom now has a metric and a documented response, which is what was missing before.

\section*{Implementation Agenda}

Adopting a feasibility-aware planning system requires the three organizational capabilities in Table~\ref{tab:implementation-agenda}. Success should be tracked across all three metric families together---forecasting (WAPE, MAE), planning (PEG rate, normalized loss), and operational (supplier rework frequency, planner coordination time)---so that feasibility gains are visible to the teams who have to approve them, not only to the team that built the model.

\begin{table}[H]
\centering
\caption{Implementation agenda for a feasibility-aware planning system.}
\label{tab:implementation-agenda}
\small
\renewcommand{\arraystretch}{0.85}
\begin{tabular}{@{}p{0.22\linewidth}p{0.72\linewidth}@{}}
\toprule
\textbf{Capability} & \textbf{Key actions} \\
\midrule
Data \& measurement & Estimate execution capacity from lead times and change cycles; log exception periods; calibrate violation cost. \\
Governance design & Define selection thresholds for WAPE, violation rate, and planning loss; stress-test before selection; align forecasting and operations on shared criteria. \\
System integration & Convert forecast-to-plan at the system boundary; embed capacity constraints in optimization; monitor violations and loss alongside WAPE. \\
\bottomrule
\end{tabular}
\end{table}


This work bridges forecasting advancement and operational execution: forecast accuracy is necessary but not sufficient for deployment. Making execution feasibility measurable and comparable across strategies lets organizations balance accuracy, stability, and operational realism, and adopt incrementally, starting with ensemble methods and progressing to DP or budgeted DP as governance maturity increases. The goal is not to replace forecasting metrics but to evaluate them where planning decisions are actually executed.

\textbf{Acknowledgements.} We thank the research collaborators and stakeholders who provided feedback on this work's operational relevance and practical applicability.

\section*{Biography}

\textbf{Yihua Xu} is an independent researcher and AI and Data professional based in San Jose, California, USA. She currently works at Accenture in the AI and Data practice, with professional interests spanning machine learning, operations research, supply chain analytics, and decision intelligence. She received her Bachelor of Science degree in Industrial Engineering from the H. Milton Stewart School of Industrial and Systems Engineering at the Georgia Institute of Technology and her Master of Business Analytics degree from the MIT Sloan School of Management at the Massachusetts Institute of Technology. Her research focuses on feasibility-aware decision frameworks for operational planning systems, including demand planning, forecast-to-plan conversion, planning stability, execution-risk evaluation, and reverse logistics. Her work aims to connect predictive models with executable operational decisions under inventory, capacity, and planning-governance constraints.

